% ============================================================
% Question Bank: ch02-03 Percent Increase and Decrease
% Topic Title: Percent Increase and Decrease
% CCSS: 7.RP.A.3
% Questions: 27 (15 MC + 6 SA + 6 Graphical)
% ============================================================

% --- Multiple Choice Questions (15) ---

\begin{practiceQuestion}{2-3-q01}{mc}
\begin{questionText}
A pair of sneakers costs $\$70$. The price increases by $20\%$. What is the new price?
\end{questionText}
\choiceA{$\$14$}
\choiceB{$\$56$}
\choiceC{$\$84$}
\choiceD{$\$90$}
\correctAnswer{C}
\explanation{Multiply by the increase multiplier: $70 \times 1.20 = \$84$. The multiplier for a $20\%$ increase is $1 + 0.20 = 1.20$.}
\end{practiceQuestion}

\begin{practiceQuestion}{2-3-q02}{mc}
\begin{questionText}
A city's population went from $40{,}000$ to $46{,}000$. What is the percent increase?
\end{questionText}
\choiceA{$6\%$}
\choiceB{$12\%$}
\choiceC{$15\%$}
\choiceD{$16\%$}
\correctAnswer{C}
\explanation{Change: $46{,}000 - 40{,}000 = 6{,}000$. Percent increase: $6{,}000 \div 40{,}000 = 0.15 = 15\%$.}
\end{practiceQuestion}

\begin{practiceQuestion}{2-3-q03}{mc}
\begin{questionText}
A laptop originally costs $\$800$. It is on sale for $15\%$ off. What is the sale price?
\end{questionText}
\choiceA{$\$120$}
\choiceB{$\$680$}
\choiceC{$\$685$}
\choiceD{$\$785$}
\correctAnswer{B}
\explanation{The decrease multiplier is $1 - 0.15 = 0.85$. Sale price: $800 \times 0.85 = \$680$.}
\end{practiceQuestion}

\begin{practiceQuestion}{2-3-q04}{mc}
\begin{questionText}
What multiplier represents a $35\%$ decrease?
\end{questionText}
\choiceA{$0.35$}
\choiceB{$0.65$}
\choiceC{$1.35$}
\choiceD{$1.65$}
\correctAnswer{B}
\explanation{A $35\%$ decrease means you keep $100\% - 35\% = 65\%$. The multiplier is $0.65$.}
\end{practiceQuestion}

\begin{practiceQuestion}{2-3-q05}{mc}
\begin{questionText}
A house's value rose from $\$150{,}000$ to $\$180{,}000$. What is the percent increase?
\end{questionText}
\choiceA{$15\%$}
\choiceB{$17\%$}
\choiceC{$20\%$}
\choiceD{$30\%$}
\correctAnswer{C}
\explanation{Change: $180{,}000 - 150{,}000 = 30{,}000$. Percent increase: $30{,}000 \div 150{,}000 = 0.20 = 20\%$.}
\end{practiceQuestion}

\begin{practiceQuestion}{2-3-q06}{mc}
\begin{questionText}
A goldfish weighed $3$ ounces and now weighs $5$ ounces. What is the percent increase in weight, rounded to the nearest whole percent?
\end{questionText}
\choiceA{$40\%$}
\choiceB{$60\%$}
\choiceC{$67\%$}
\choiceD{$167\%$}
\correctAnswer{C}
\explanation{Change: $5 - 3 = 2$. Percent increase: $2 \div 3 \approx 0.667 = 66.7\% \approx 67\%$. Always divide by the original.}
\end{practiceQuestion}

\begin{practiceQuestion}{2-3-q07}{mc}
\begin{questionText}
What multiplier represents a $12\%$ increase?
\end{questionText}
\choiceA{$0.12$}
\choiceB{$0.88$}
\choiceC{$1.12$}
\choiceD{$12$}
\correctAnswer{C}
\explanation{A $12\%$ increase: $1 + 0.12 = 1.12$.}
\end{practiceQuestion}

\begin{practiceQuestion}{2-3-q08}{mc}
\begin{questionText}
A smartphone cost $\$480$ last year. This year the price dropped to $\$408$. What is the percent decrease?
\end{questionText}
\choiceA{$12\%$}
\choiceB{$15\%$}
\choiceC{$18\%$}
\choiceD{$72\%$}
\correctAnswer{B}
\explanation{Change: $480 - 408 = 72$. Percent decrease: $72 \div 480 = 0.15 = 15\%$.}
\end{practiceQuestion}

\begin{practiceQuestion}{2-3-q09}{mc}
\begin{questionText}
After a $40\%$ increase, a gym membership costs $\$84$ per month. What was the original price?
\end{questionText}
\choiceA{$\$50.40$}
\choiceB{$\$56$}
\choiceC{$\$60$}
\choiceD{$\$62$}
\correctAnswer{C}
\explanation{After a $40\%$ increase the new price is $1.40$ times the original. Divide: $84 \div 1.40 = \$60$.}
\end{practiceQuestion}

\begin{practiceQuestion}{2-3-q10}{mc}
\begin{questionText}
A farmer had $250$ chickens. After one year he has $200$ chickens. What is the percent decrease?
\end{questionText}
\choiceA{$15\%$}
\choiceB{$20\%$}
\choiceC{$25\%$}
\choiceD{$50\%$}
\correctAnswer{B}
\explanation{Change: $250 - 200 = 50$. Percent decrease: $50 \div 250 = 0.20 = 20\%$.}
\end{practiceQuestion}

\begin{practiceQuestion}{2-3-q11}{mc}
\begin{questionText}
Darius says: ``A $20\%$ increase followed by a $20\%$ decrease returns the price to its original value.'' Is he correct?
\end{questionText}
\choiceA{Yes, the changes cancel out}
\choiceB{No, the final price is less than the original}
\choiceC{No, the final price is more than the original}
\choiceD{It depends on the starting price}
\correctAnswer{B}
\explanation{Use $\$100$: increase $20\%$ gives $\$120$; decrease $20\%$ of $\$120$ gives $120 \times 0.80 = \$96$. The final price is less because the decrease is applied to a larger number.}
\end{practiceQuestion}

\begin{practiceQuestion}{2-3-q12}{mc}
\begin{questionText}
A hotel room costs $\$120$ per night during the off-season. During peak season the price is $\$156$. What is the percent increase?
\end{questionText}
\choiceA{$20\%$}
\choiceB{$25\%$}
\choiceC{$30\%$}
\choiceD{$36\%$}
\correctAnswer{C}
\explanation{Change: $156 - 120 = 36$. Percent increase: $36 \div 120 = 0.30 = 30\%$.}
\end{practiceQuestion}

\begin{practiceQuestion}{2-3-q13}{mc}
\begin{questionText}
Which of the following changes results in the greatest percent decrease?
\end{questionText}
\choiceA{From $\$100$ to $\$75$}
\choiceB{From $\$200$ to $\$160$}
\choiceC{From $\$80$ to $\$56$}
\choiceD{From $\$50$ to $\$38$}
\correctAnswer{C}
\explanation{A: $25\%$ decrease. B: $20\%$ decrease. C: $(80-56) \div 80 = 30\%$ decrease. D: $(50-38) \div 50 = 24\%$ decrease. C is greatest.}
\end{practiceQuestion}

\begin{practiceQuestion}{2-3-q14}{mc}
\begin{questionText}
After a $25\%$ discount, a coat costs $\$63$. What was the original price?
\end{questionText}
\choiceA{$\$78.75$}
\choiceB{$\$80$}
\choiceC{$\$84$}
\choiceD{$\$88$}
\correctAnswer{C}
\explanation{After $25\%$ off you pay $75\%$: $63 \div 0.75 = \$84$.}
\end{practiceQuestion}

\begin{practiceQuestion}{2-3-q15}{mc}
\begin{questionText}
A sunflower was $18$ inches tall and grew to $27$ inches. What is the percent increase?
\end{questionText}
\choiceA{$33\%$}
\choiceB{$40\%$}
\choiceC{$50\%$}
\choiceD{$67\%$}
\correctAnswer{C}
\explanation{Change: $27 - 18 = 9$. Percent increase: $9 \div 18 = 0.50 = 50\%$.}
\end{practiceQuestion}

% --- Short Answer Questions (6) ---

\begin{practiceQuestion}{2-3-q16}{sa}
\begin{questionText}
A bicycle costs $\$360$. It is marked down $35\%$. What is the sale price?
\end{questionText}
\correctAnswer{$\$234$}
\explanation{Multiply by $1 - 0.35 = 0.65$: $360 \times 0.65 = \$234$.}
\end{practiceQuestion}

\begin{practiceQuestion}{2-3-q17}{sa}
\begin{questionText}
The temperature went from $80^\circ$F to $68^\circ$F. What is the percent decrease?
\end{questionText}
\correctAnswer{$15\%$}
\explanation{Change: $80 - 68 = 12$. Percent decrease: $12 \div 80 = 0.15 = 15\%$.}
\end{practiceQuestion}

\begin{practiceQuestion}{2-3-q18}{sa}
\begin{questionText}
What multiplier represents a $45\%$ increase?
\end{questionText}
\correctAnswer{$1.45$}
\explanation{Add the percent as a decimal to $1$: $1 + 0.45 = 1.45$.}
\end{practiceQuestion}

\begin{practiceQuestion}{2-3-q19}{sa}
\begin{questionText}
A store's monthly revenue increased from $\$8{,}000$ to $\$10{,}400$. What is the percent increase?
\end{questionText}
\correctAnswer{$30\%$}
\explanation{Change: $10{,}400 - 8{,}000 = 2{,}400$. Percent increase: $2{,}400 \div 8{,}000 = 0.30 = 30\%$.}
\end{practiceQuestion}

\begin{practiceQuestion}{2-3-q20}{sa}
\begin{questionText}
A video game was $\$50$. After a price increase, it costs $\$62.50$. What was the percent increase?
\end{questionText}
\correctAnswer{$25\%$}
\explanation{Change: $62.50 - 50 = 12.50$. Percent increase: $12.50 \div 50 = 0.25 = 25\%$.}
\end{practiceQuestion}

\begin{practiceQuestion}{2-3-q21}{sa}
\begin{questionText}
After a $30\%$ decrease, a TV costs $\$350$. What was the original price?
\end{questionText}
\correctAnswer{$\$500$}
\explanation{After a $30\%$ decrease you pay $70\%$: $350 \div 0.70 = \$500$.}
\end{practiceQuestion}

% --- Graphical Questions (3 GMC + 3 GSA) ---

\begin{practiceQuestion}{2-3-q22}{gmc}
\begin{questionText}
The bar graph shows the population of a town over two years. What is the percent increase from Year 1 to Year 2?

\begin{center}
\begin{tikzpicture}[scale=0.5]
  \draw[thick,->] (0,0) -- (0,7) node[above]{\small Population};
  \draw[thick,->] (0,0) -- (6,0);
  \fill[funBlue!60] (1,0) rectangle (2.5,5);
  \fill[funGreen!60] (3.5,0) rectangle (5,6);
  \node[below,font=\small] at (1.75,0) {Year 1};
  \node[below,font=\small] at (4.25,0) {Year 2};
  \foreach \y/\lab in {1/2000,2/4000,3/6000,4/8000,5/10000,6/12000} {
    \draw (-0.1,\y) -- (0.1,\y) node[left,font=\small,xshift=-2pt] {\lab};
  }
  \node[above,font=\small] at (1.75,5) {$10{,}000$};
  \node[above,font=\small] at (4.25,6) {$12{,}000$};
\end{tikzpicture}
\end{center}
\end{questionText}
\choiceA{$2\%$}
\choiceB{$12\%$}
\choiceC{$17\%$}
\choiceD{$20\%$}
\correctAnswer{D}
\explanation{Change: $12{,}000 - 10{,}000 = 2{,}000$. Percent increase: $2{,}000 \div 10{,}000 = 0.20 = 20\%$.}
\end{practiceQuestion}

\begin{practiceQuestion}{2-3-q23}{gmc}
\begin{questionText}
The table shows the price of a skateboard during two sales. Which sale offers the greater percent decrease from the original price?

\begin{center}
\begin{tabular}{|l|c|c|}
\hline
 & \textbf{Original Price} & \textbf{Sale Price} \\
\hline
Sale A & $\$80$ & $\$60$ \\
\hline
Sale B & $\$120$ & $\$84$ \\
\hline
\end{tabular}
\end{center}
\end{questionText}
\choiceA{Sale A ($20\%$ decrease)}
\choiceB{Sale A ($25\%$ decrease)}
\choiceC{Sale B ($30\%$ decrease)}
\choiceD{Both are the same percent decrease}
\correctAnswer{C}
\explanation{Sale A: $(80-60) \div 80 = 25\%$. Sale B: $(120-84) \div 120 = 30\%$. Sale B has the greater percent decrease.}
\end{practiceQuestion}

\begin{practiceQuestion}{2-3-q24}{gmc}
\begin{questionText}
The number line below shows an original value and a new value. What is the percent decrease?

\begin{center}
\begin{tikzpicture}[scale=0.09]
  \draw[thick,->] (0,0) -- (105,0);
  \foreach \x in {0,10,...,100} {
    \draw (\x,1.5) -- (\x,-1.5) node[below,font=\small] {\x};
  }
  \fill[funRed] (90,0) circle (3) node[above=10pt,font=\small] {Original};
  \fill[funBlue] (72,0) circle (3) node[above=10pt,font=\small] {New};
  \draw[thick,->,funBlue] (88,12) -- (74,12);
\end{tikzpicture}
\end{center}
\end{questionText}
\choiceA{$18\%$}
\choiceB{$20\%$}
\choiceC{$25\%$}
\choiceD{$28\%$}
\correctAnswer{B}
\explanation{Change: $90 - 72 = 18$. Percent decrease: $18 \div 90 = 0.20 = 20\%$. Always divide by the original.}
\end{practiceQuestion}

\begin{practiceQuestion}{2-3-q25}{gsa}
\begin{questionText}
The bar graph shows the number of library visitors on two days. What is the percent increase from Monday to Tuesday?

\begin{center}
\begin{tikzpicture}[scale=0.5]
  \draw[thick,->] (0,0) -- (0,7) node[above]{\small Visitors};
  \draw[thick,->] (0,0) -- (6,0);
  \fill[funOrange!60] (1,0) rectangle (2.5,4);
  \fill[funBlue!60] (3.5,0) rectangle (5,5);
  \node[below,font=\small] at (1.75,0) {Mon};
  \node[below,font=\small] at (4.25,0) {Tue};
  \foreach \y/\lab in {1/30,2/60,3/90,4/120,5/150,6/180} {
    \draw (-0.1,\y) -- (0.1,\y) node[left,font=\small,xshift=-2pt] {\lab};
  }
  \node[above,font=\small] at (1.75,4) {$120$};
  \node[above,font=\small] at (4.25,5) {$150$};
\end{tikzpicture}
\end{center}
\end{questionText}
\correctAnswer{$25\%$}
\explanation{Change: $150 - 120 = 30$. Percent increase: $30 \div 120 = 0.25 = 25\%$.}
\end{practiceQuestion}

\begin{practiceQuestion}{2-3-q26}{gsa}
\begin{questionText}
The table shows how many members a club had at the start and end of the school year. Find the percent decrease.

\begin{center}
\begin{tabular}{|l|c|}
\hline
\textbf{Time} & \textbf{Members} \\
\hline
Start of year & $60$ \\
\hline
End of year & $45$ \\
\hline
\end{tabular}
\end{center}
\end{questionText}
\correctAnswer{$25\%$}
\explanation{Change: $60 - 45 = 15$. Percent decrease: $15 \div 60 = 0.25 = 25\%$.}
\end{practiceQuestion}

\begin{practiceQuestion}{2-3-q27}{gsa}
\begin{questionText}
The bar model below shows a price \textbf{before} and \textbf{after} a percent decrease. What was the percent decrease?

\begin{center}
\begin{tikzpicture}[scale=0.6]
  % Original bar
  \draw[thick] (0,1.2) rectangle (10,2);
  \fill[funBlue!30] (0,1.2) rectangle (10,2);
  \node[above,font=\small] at (5,2) {Original: $\$160$};
  % New bar
  \draw[thick] (0,0) rectangle (7,0.8);
  \fill[funGreen!30] (0,0) rectangle (7,0.8);
  \node[below,font=\small] at (3.5,0) {Sale: $\$112$};
  % Removed portion
  \draw[thick,dashed] (7,0) rectangle (10,0.8);
  \node[font=\small] at (8.5,0.4) {?};
\end{tikzpicture}
\end{center}
\end{questionText}
\correctAnswer{$30\%$}
\explanation{Change: $160 - 112 = 48$. Percent decrease: $48 \div 160 = 0.30 = 30\%$.}
\end{practiceQuestion}
